\Needspace{6\baselineskip}
\begin{center}\hypertarget{native-254-article181}{}\pdfbookmark[2]{181.}{article-181}181.\end{center}
Gibt es also nur eine Unbekannte $p$ und hat man zur Bestimmung dieser Unbekannten die Werte der Funktionen $ap+n$, $a'p+n'$, $a''p+n''$ usw. durch gleich genaue Beobachtungen jeweils als $=M$, $M'$, $M''$ usw. ermittelt, so ist der wahrscheinlichste Wert von $p$
\[
=\frac{am+a'm'+a''m''+\text{usw.}}{aa+a'a'+a''a''+\text{usw.}}=A,
\]
wobei $m$, $m'$, $m''$ usw. jeweils für $M-n$, $M'-n'$, $M''-n''$ usw. stehen.

Um abzuschätzen, welcher Grad an Präzision bei diesem Wert zu erwarten ist, nehmen wir an, dass die Wahrscheinlichkeit eines Beobachtungsfehlers $\Delta$ durch $\frac{h}{\BareRoot{\pi}}e^{-hh\Delta\Delta}$ gegeben ist. Die Wahrscheinlichkeit dafür, dass der wahre Wert von $p$ $=A+p'$ ist, ist daher proportional zur Funktion
\[
e^{-hh((ap-m)^2+(a'p-m')^2+(a''p-m''^2+\text{usw.})},
\]
wenn man $p$ durch $A+p'$ ersetzt. Der Exponent dieser Funktion lässt sich auf die Form $-hh(aa+a'a'+a''a''+\text{usw.})(pp-2pA+B)$ bringen, wobei $B$ von $p$ unabhängig ist. Die Funktion selbst ist daher proportional zu
\[
e^{-hh(aa+a'a'+a''a''+\text{usw.})p'p'}.
\]
Dem Wert $A$ ist also derselbe Grad an Präzision zuzuschreiben, als hätte man ihn durch eine unmittelbare Beobachtung gefunden, deren Präzision sich zu der der ursprünglichen Beobachtungen wie $h\BareRoot{(aa+a'a'+a''a''+\text{usw.})}$ zu $h$ oder wie $\BareRoot{(aa+a'a'+a''a''+\text{usw.})}$ zu $1$ verhält.
